% 第5章图：按次序(1,4,2,3)对四变量例子消元时函数集合的逐步变化（因子图）。
% psi1(x2,x3)=(0,0,1,2), b1=(1,0,1,1)；psi4(x2,x3)=(1,0,0,1), b4=(1,0,1,0)；psi2'(x3)=(1,0), b2=(0,0)；常数0, b3=1。
% 逆序回溯：x3=1, x2=0, x4=0, x1=0。
\begin{tikzpicture}[x=1mm,y=1mm,>=Stealth,
  v/.style={circle,draw=FigureStroke,fill=FigureBlueFill,line width=.7pt,minimum size=5.2mm,inner sep=0pt},
  cur/.style={v,fill=FigureYellowFill,draw=FigureYellow},
  ghost/.style={circle,draw=FigureGrid,dashed,line width=.5pt,minimum size=5.2mm,inner sep=0pt,text=FigureGrid},
  f/.style={rectangle,draw=FigureStroke,fill=white,line width=.7pt,minimum width=6.4mm,minimum height=4.6mm,inner sep=.5pt,font=\figurefont\footnotesize},
  fn/.style={f,fill=FigureGreenFill,draw=FigureGreen},
  e/.style={FigureStroke,line width=.6pt},
  tt/.style={font=\figurefont\footnotesize,text=FigureInk},
  val/.style={font=\figurefont\footnotesize,text=FigureInk,anchor=west},
  note/.style={font=\FigureNoteFont\fontsize{7.5}{9}\selectfont,text=FigureMuted}]
\def\P{30.5}
% ---------- panel 0: initial ----------
\begin{scope}
 \node[tt] at (9,21) {初始};
 \node[cur] (a1) at (0,0) {$1$}; \node[v] (a2) at (6,0) {$2$}; \node[v] (a3) at (12,0) {$3$}; \node[v] (a4) at (18,0) {$4$};
 \node[f] (fA) at (3.5,11) {$\phi_A$}; \node[f] (fB) at (14.5,11) {$\phi_B$};
 \draw[e] (fA)--(a1) (fA)--(a2) (fA)--(a3) (fB)--(a2) (fB)--(a3) (fB)--(a4);
 \node[note,align=center] at (9,-8.5) {黄色变量为\\下一步消去者};
\end{scope}
\draw[->,FigureStroke,line width=.7pt] (22.5,5)--(27.5,5);
% ---------- panel 1: after x1 ----------
\begin{scope}[xshift=\P mm]
 \node[tt] at (9,21) {消去$x_1$后};
 \node[ghost] (b1) at (0,0) {$1$}; \node[v] (b2) at (6,0) {$2$}; \node[v] (b3) at (12,0) {$3$}; \node[cur] (b4) at (18,0) {$4$};
 \node[fn] (g1) at (5,11) {$\psi_1$}; \node[f] (gB) at (14.5,11) {$\phi_B$};
 \draw[e] (g1)--(b2) (g1)--(b3) (gB)--(b2) (gB)--(b3) (gB)--(b4);
 \node[note,align=center] at (9,-8.5) {$\psi_1=(0,0,1,2)$\\$b_1=(1,0,1,1)$};
\end{scope}
\draw[->,FigureStroke,line width=.7pt] (\P+22.5,5)--(\P+27.5,5);
% ---------- panel 2: after x4 ----------
\begin{scope}[xshift=2*\P mm]
 \node[tt] at (9,21) {消去$x_4$后};
 \node[ghost] (c1) at (0,0) {$1$}; \node[cur] (c2) at (6,0) {$2$}; \node[v] (c3) at (12,0) {$3$}; \node[ghost] (c4) at (18,0) {$4$};
 \node[f] (h1) at (3.5,11) {$\psi_1$}; \node[fn] (h4) at (14.5,11) {$\psi_4$};
 \draw[e] (h1)--(c2) (h1)--(c3) (h4)--(c2) (h4)--(c3);
 \node[note,align=center] at (9,-8.5) {$\psi_4=(1,0,0,1)$\\$b_4=(1,0,1,0)$};
\end{scope}
\draw[->,FigureStroke,line width=.7pt] (2*\P+22.5,5)--(2*\P+27.5,5);
% ---------- panel 3: after x2 ----------
\begin{scope}[xshift=3*\P mm]
 \node[tt] at (9,21) {消去$x_2$后};
 \node[ghost] at (0,0) {$1$}; \node[ghost] at (6,0) {$2$}; \node[cur] (d3) at (12,0) {$3$}; \node[ghost] at (18,0) {$4$};
 \node[fn] (k2) at (12,11) {$\psi_2'$};
 \draw[e] (k2)--(d3);
 \node[note,align=center] at (9,-8.5) {$\psi_2'=(1,0)$\\$b_2=(0,0)$};
\end{scope}
\draw[->,FigureStroke,line width=.7pt] (3*\P+22.5,5)--(3*\P+27.5,5);
% ---------- panel 4: after x3 ----------
\begin{scope}[xshift=4*\P mm]
 \node[tt] at (9,21) {消去$x_3$后};
 \node[ghost] at (0,0) {$1$}; \node[ghost] at (6,0) {$2$}; \node[ghost] at (12,0) {$3$}; \node[ghost] at (18,0) {$4$};
 \node[fn] at (9,11) {$0$};
 \node[note,align=center] at (9,-8.5) {最小成本$0$\\$b_3=1$};
\end{scope}
% ---------- backtracking ----------
\draw[FigureGrid,line width=.5pt] (-3,-16)--(4*\P+21,-16);
\node[note,anchor=west] at (-3,-20.5) {逆序回溯};
\node[tt,anchor=west] at (16,-20.5) {$x_3=b_3=1\ \to\ x_2=b_2(1)=0\ \to\ x_4=b_4(0,1)=0\ \to\ x_1=b_1(0,1)=0$};
\end{tikzpicture}
